\chapter*{Introducere} 
\addcontentsline{toc}{chapter}{Introducere}
În ciuda faptului că dinamica populației a atras din punct de vedere istoric interesul multor oameni de știință, suntem încă departe de a înțelege cum și de ce populațiile fluctuează în ecosistemele naturale. Motivul principal este că populațiile sunt sisteme complexe, cu proprietăți emergente care sunt imposibil de determinat prin suma părților lor.\cite{ChallengesInPopulationDynamics}\\
Intr-un mediu virtual, am putea izola si controla variabilele din mediu pentru a studia anumiti factori specifici care influenteaza aceasta dinamica. Specific, vom simula si studia interactiunile prada-pradator, impactul lor reflectându-se in dinamica populatiei in urmatoarele moduri:
\begin{itemize}
    \item \textbf{Reglarea populatiei} - prădarea reprezinta un mecanism regulator atat pentru populatia prazii cat si a pradatorilor, ce previne prada din a ajunge la densitati prea mari, iar densitatea prădătorilor e controlata de cea a prazii.
    \item \textbf{Dinamica ciclica a populatiei} - interactiunile prada-pradator pot rezulta in ciclicitate a populatiilor, unde numarul prazii si a pradatorilor pot sa prezinte fluctuatii in timp in mod regulat. In studierea impactului pe care l-ar avea incalzirea globala asupra dinamicii sistemelor ecologice sa observat chiar ca durata acestor cicluri se micsoreaza trecand de la temperaturi calde la cele mai reci\cite{TemperatureCycle}.
    \item \textbf{Dinamici evolutive} - presiunea impusa de mecanismul de pradare poate duce la adaptari evolutive. In mediul natural prăzile pot dezvolta mecanisme defensive(coarne, camuflaj, viteza) iar pradatorii deasemenea pot suferi schimbari fizice dar si in strategiile de vanatoare.
\end{itemize}

În aceasta lucrare vom dezvolta un mediu virtual cu o populație de prada și una de prădători, cu scopul de a analiza transformarile suferite in populatii dupa criteriile de mai sus. Mediul va fi dezvoltat in motorul de jocuri Unity, si il vom descrie in capitolul \ref{MediulVirtual}.El va include 4 pereti pozitionati pentru a crea o arie inchisa. In cadrul acelei arii vor fi generati in mod constant peleti de energie care vor fi consumati de prada, va fi generata populatia de prada si populatia de pradatori. Evolutia organismelor din fiecare populatie se va preta pe dezvoltarea comportamentelor, nu a trasaturilor fizice. Pentru a directiona aceasta evolutie, vom folosi un algoritm din clasa algoritmilor evolutivi.

De-a lungul dezvoltarii familiei de algoritmi evolutivi, apar anumite etape cheie in care se dezvolta ramuri noi. Enumeram cateva din ramuri in urmatorul mod:
\begin{itemize}
    \item \textbf{Algoritmii Genetici}(GA) - arie de studiu în care pionier este John Holland prin cartea sa "Adaptation in Natural and Artificial Systems", publicată în 1975\cite{GA}. In ea el introduce utilizarea unei populatii de indivizi, a operatorilor genetici precum mutatie si crossover, si a mecanismelor de selectie inspirate de selectia naturala.
    \item \textbf{Strategii Evolutive}(ES) - dezvoltarea acestei ramuri poate fi atribuita lui Rechenberg și Hans-Paul Schwefel care sau ocupat de studiul ei în anii '60 și '70.
    In 1973 Rechenberg publica cartea "Evolutionsstrategie"\cite{ES}, in care documenteaza cercetarea sa pe probleme de optimizare continua.
    \item \textbf{NeuroEvolutie} - aceasta clasa de algoritmi cauta prin spatiul de ponderi si structuri o retea neuronala cu performante bune pentru a anumita sarcina. Practic algoritmii de neuroevolutie permit gasirea unui comportament care faciliteaza rezolvarea problemelor de \textbf{reinforcement learning} complexe\cite{SymbioticEvolution}. De obicei, in cadrul algoritmilor NE, se alege o tipologie(structura de noduri si conexiuni) pentru retele ce urmeaza sa treaca prin procesul de evolutie, iar schimbarile au loc doar pe ponderile conexiunilor. \textbf{NEAT} este o abordare inedita a neuroevolutiei pe parcursul careia se modifica atat tipologia retelei cat si ponderile conexiunilor. Astfel este scutit efortul uman de a crea designul unei retele. La inceptia acestei metode a stat Kenneth O. Stanley, cu articolul s-au din 2002 "Evolving Neural Networks through Augmenting Topologies"\cite{NEAT}
\end{itemize}

In cadrul acestei teze, comportamentele fiecarui individ din populatie, adica actiunile pe care le fac la fiecare pas de timp vor fi dictate de cate o retea neuronala. Ca algoritmi de neuroevolutie vom folosi \textbf{NEAT} si \textbf{rtNEAT}. Acesti algoritmi vor fi descrisi fiecare in capitolele \ref{NEAT} si \ref{rtNEAT}.In cadrul unui experiment in care vom include doar speciile de prada, in care excludem analiza populatiei ca numar de indivizi, vom identifica care dintre algoritmi determina o progresie mai rapida in materie de acumulare de energie in medie si cata energie acumuleaza cel mai bun membru al populatiei. Un experiment similar a fost descris in lucrarea lui T. Gelles si M. Sanchez\cite{ants} in care incearca sa evolueze comportamentul furnicilor de a plasa feromoni ca un mod de a comunica, si demonstreaza ca \textbf{rtNEAT} are mai mult succes in a evolua acest comportament decat \textbf{NEAT}.\\
Pentru alte experimente in materie de interactiune prada-pradator putem face referinta la lucrarea "An evolutionary approach for interactive computer games"\cite{DeadEndGame}. In cadrul ei autorii proiecteaza un joc numit "Dead End", in care jucatorul are scopul de a ajunge la iesirea din partea opusa a ecranului de unde este el generat, iar o echipa de 8 "Dogs", ai caror miscari sunt dictate de cate o retea neuronala, au scopul de al prinde inainte ca acesta sa ajunga la iesire, deci trebuiau sa dezvolte comportamente de urmarire si de cooperare. Pentru evolutia cainilor au fost folositi 2 algoritmi de invatare proprii, unul offline si unul online. Ambii aduceau schimbări doar ponderilor conexiunilor din retea, nu si a topologiei retelelor. Ei au identificat ca agentii antrenati in modul online prezinta o adapabilitate mare la schimbarea strategiilor de joc a jucatorilor, deci demonstreaza eficacitatea metodei impotriva jucatorilor umani.\\
Datorita concluziilor acestor 2 lucrari, vom sustine ipoteza ca \textbf{rtNEAT}, care este un algoritm de invatare online, va determina o evolutie mai rapida a indivizilor, in termeni de energie acumulata. Deasemenea, doar\textbf{rtNEAT} ne permite sa facem o analiza a populatiilor pe timp indelungat, deoarece \textbf{NEAT} reseteaza populatile la fiecare generatie.
Daca identificam ca \textbf{NEAT}, care este un algoritm offline, produce o invatare mai rapida, il vom folosi pentru a evolua comportamentele indivizilor, iar retelele antrenate astfel le vom integra in \textbf{rtNEAT} pentru a face analiza dinamicii populatiei.
In capitolul \ref{DetaliiImplementare} vom detalia implementarea mediului, a algoritmilor, a retelelor, si a senzorilor folositi. In capitolul \ref{SetariExperimentale} vom prezenta experimentele pe care vrem sa le conducem si setarile pentru algoritmi in cadrul experimentelor. In final, \ref{Rezultate} prezentam rezultatele si anlizam care sunt factorii care le-au influentat.
